%قبل از شروع مسئله ترجیحاً چیزی ننویسید. اما اگر خواستید بنویسید، در قالب لیستی (فقرات و شمارش) بنویسید.
%برای اضافه کردن نمره، بعدی از شروع‌مسئله، [x] را اضافه کنید و به جای x نمره مورد نظر را قرار دهید. برای زیرمسئله‌ها هم به همین صورت است.
%\شروع‌مسئله[10]
\شروع‌مسئله [15]
%\متن‌سیاه{سوال عملی اول:}\\
%\متن‌سیاه{دراپ‌تست}\\

%\متن‌سیاه{ورودی:}\\
%در خط اول ورودی، ابتدا k و سپس m داده می‌شود.\\
%\متن‌سیاه{خروجی:}\\
%در تنها سطر خروجی، تعداد آیفون‌های مورد نیاز می‌آید. اگر در هیچ صورتی چنین کاری ممکن نبود ۱- چاپ می‌شود.\\
%\متن‌سیاه{محدودیت‌ها:}\\
%$1\leq k\leq 10^4$

%$m\leq10^9$
%\\
%\textbf{\متن‌سیاه{ورودی و خروجی نمونه}\\
%% تست اول
%\begin{latin}
% \begin{tabularx}{\textwidth}{|X | X|}
% \hline
% \textbf{stdin} &
% \textbf{stdout}\\
% \hline
% 100 100&
% 1\\
% \hline
% 
% \end{tabularx}
%\end{latin}
%
%% تست دوم
%
%\begin{latin}
% \begin{tabularx}{\textwidth}{|X | X|}
% \hline
% \textbf{stdin} &
% \textbf{stdout}\\
% \hline
% 340 20&
% 3\\
% \hline
% 
% \end{tabularx}
%\end{latin}}
% صورت اصلی مسئله را در این‌جا بنویسید.

شنگدباو در گرافی $n$ رأسی و $m$ یالی وزن‌دار و جهت‌دار گیر افتاده‌است! او در ابتدا در رأس $s$ قرار دارد و برای خروج باید به رأس $t$ برود که در آن دری به خارج از گراف وجود دارد. نکته‌ای که وجود دارد این است که در مورد نظر در ابتدا باز است اما با گذشتن هر ثانیه در تغییر وضعیت می‌دهد یعنی در ثانیه اول باز است و در ثانیه دوم بسته و در ثانیه سوم باز و ... 

شنگدباو باید در هر لحظه حرکت کند و مجاز نیست سر جای خود بایستد. همچنین یال‌ها یکطرفه هستند لذا روی هر یال صرفاً به یک جهت می‌تواند حرکت کند. وزن نوشته‌شده روی هر یال برابر طول آن یال است یعنی اگر وزن یالی برابر $c$ باشد، شنگدباو برای رسیدن از ابتدا به انتهای آن باید $c$ ثانیه بدود!

به او کمک کنید تا کمترین زمانی که طول می‌کشد تا او از شر این گراف خلاص شود را بیابد.

\noindent
\متن‌سیاه{ورودی:}\\
در خط اول ورودی دو عدد $n$ و $m$ به ترتیب داده‌می‌شود که به معنای تعداد رأس‌ها و یال‌های گراف مورد نظر است. سپس در خط بعدی شماره دو رأس $s$ و $t$ به ترتیب داده می‌شود.

در $m$ خط بعدی در هر خط سه عدد $u$ و $v$ و $c$ می‌آیند که بیانگر رأس ابتدایی و انتهایی و وزن یال متناظر است.

\noindent
\متن‌سیاه{خروجی:}\\
اگر شنگدباو نمی‌توانست از گراف خارج شود $-1$ خروجی دهید و در غیر این‌صورت کمترین تعداد ثانیه‌هاییکه طول می‌کشد تا او از گراف خارج شود را خروجی دهید.

\متن‌سیاه{محدودیت‌ها:}\\
$1 \le n, m \le 10^5$

وزن یال‌ها نیز حداکثر $10^9$ می‌باشد.

\متن‌سیاه{ورودی و خروجی نمونه:}\\
%% تست اول
\begin{latin}
\begin{tabularx}{\textwidth}{|X | X|}
	\hline
	\textbf{stdin} &
	\textbf{stdout}\\
	\hline
	3 6
	
	1 2
	
	1 2 1
	
	2 1 1
	
	1 3 1001
	
	3 1 1001
	
	3 2 1001
	
	2 3 1001&
	2002\\
	\hline	
\end{tabularx}
\end{latin}

\solution{
}

%%%%%%%%%%%%%%%%%%%%%%%%%%%%%%%%%%%%%%%%%%%%%%%%%%%%%
% برخی دستورات مفید برای تمرین‌ها. %
%%%%%%%%%%%%%%%%%%%%%%%%%%%%%%%%%%%%%%%%%%%%%%%%%%%%%

% فرمول‌های نمونه:
%$f(x) = \Theta(x)$
%$T(x, y) \leq \log(x) + 2^y + \binom{2n}{n}$
% {\tt code\_function}


% لیست موارد بدون شماره‌گذاری:
%\شروع{فقرات}
% \فقره مورد اول 
% \شروع{فقرات}
% \فقره مورد اول داخلی 
% \شروع{فقرات}
% \فقره اولین مورد خیلی داخلی 
% \فقره دومین مورد خیلی داخلی 
% \پایان{فقرات}
% \فقره مورد دوم داخلی 
% \پایان{فقرات}
% \فقره مورد دوم 
%\پایان{فقرات}

% لیست موارد شماره‌گذاری شده:
%\شروع{شمارش}
% \فقره یک مورد 
% \فقره یک مورد دیگر 
% \فقره باز هم یک مورد دیگر
%\پایان{شمارش}

% برای نوشتن یک فرمول چند خطی که یک شماره بخوره. برای حالت‌های دیگر به سایت زیر مراجعه کنید:
% https://www.overleaf.com/learn/latex/Aligning_equations_with_amsmath
%\begin{align} 
% \begin{split}
% (x+y)^3 &= (x+y)^2(x+y) \\
% &= (x^2+2xy+y^2)(x+y) \\
% &= (x^3+2x^2y+xy^2) + (x^2y+2xy^2+y^3) \\
% &= x^3+3x^2y+3xy^2+y^3
% \end{split} 
%\end{align}

% نوشتن ماتریس و آرایه دوبعدی
%\begin{align}
% A = 
% \begin{bmatrix}
% A_{11} & A_{21} \\
% A_{21} & A_{22}
% \end{bmatrix}
%\end{align}

% اضافه کردن تصویر با شرح. برای اضافه کردن تصویر بدون شرح، فقط دستور درج‌تصویر را استفاده کنید.
%\شروع{شکل}[h!]
%\تنظیم‌ازوسط
% \درج‌تصویر[پهنا=0.5\پهنای‌متن]{img.jpg}
%\شرح{شرح تصویر}
%\برچسب{تصویر:۱}
%\پایان{شکل}
%\پایان{زیرمسائل}

%کد زیر را ببینید.
%%\lstinputlisting{source_filename.py}
%\begin{latin}
%\begin{lstlisting}
%# Import the modules
%import sys
%import random
%
%ans = True
%
%while ans:
%question = raw_input("Ask the magic 8 ball a question: (press enter to quit) ")
%
%answers = random.randint(1,8)
%
%if question == "":
%sys.exit()
%
%elif answers == 1:
%print "It is certain"
%
%elif answers == 2:
%print "Outlook good"
%
%elif answers == 3:
%print "You may rely on it"
%
%elif answers == 4:
%print "Ask again later"
%
%elif answers == 5:
%print "Concentrate and ask again"
%
%elif answers == 6:
%print "Reply hazy, try again"
%
%elif answers == 7:
%print "My reply is no"
%
%elif answers == 8:
%print "My sources say no"
%\end{lstlisting}
%\end{latin}